\section{Implementation in ProvSQL}\label{sec:implementation}

Theorem~\ref{thm:rewriting} and Corollary~\ref{cor:pqe} pave the way for
a practical implementation of provenance management and probabilistic
query evaluation in a SQL DBMS. Note that the data model follows as
closely as possible the practical behavior of SQL engines, instead of the
theoretical model of annotated relations used in most of the literature:
multiset semantics, explicit duplicate elimination using
\mintinline{sql}{DISTINCT} or \mintinline{sql}{GROUP BY},
ordered attributes within relations.

We now explain how ProvSQL is implemented. ProvSQL uses PostgreSQL's
extension mechanism to change the behavior of the DBMS. This extension mechanism
allows to provide user-defined functions (UDFs), which can be written
either in PL/pgSQL, the application language of PostgreSQL, or in~C, to
implement provenance- and probability-related functionalities in~C. The
extension
mechanism also allows to run \emph{hooks} through a C interface at
different key phases of query evaluation. ProvSQL uses a combination of
SQL, PL/pgSQL, C, and C++ code interfacing with the C code, to implement
its behavior.

\subsection{Annotated relations}
\label{sec:impl_relations}

As in Section~\ref{sec:semirings}, ProvSQL adds to every relation for
which provenance needs to be tracked (which can be specified with a UDF)
an extra \emph{provsql} attribute
to store the annotation. Instead of choosing a specific algebraic
structure for the annotation ahead of time, this annotation is a generic
\emph{universally unique identifier
	(UUID)}~\cite{DBLP:journals/rfc/rfc9562}. We explain further on how these
generic annotations can be specialized to specific (m)-semirings.
Annotations of \emph{base tuples} are randomly generated using the UUIDv4
standard and can be interpreted as abstract tuple identifiers; the
128-bit address space makes the collision probability vanishingly
small.
Annotations are built from these base annotations by using the
(m-)semiring operations; to store resulting annotations in the same
attribute, they are also UUIDs, this time computed following UUIDv5
by computing a \mbox{SHA-1} hash formed from a fixed namespace, and a
normalized description of the operator and its UUID operands. This
guarantees, in particular, that annotations are generated in a
deterministic manner: the same query over the same data will result in
the exact same result, annotations included.

Retaining opaque UUIDs is not enough: one also needs to be able to
efficiently determine how they relate to each other. This is done by
storing a \emph{provenance circuit} (a compact DAG representation of
provenance expressions, first introduced in~\cite{deutch2014circuits})
where each UUIDv4 points to a leaf gate representing inputs of the
circuit and each UUIDv5 to an inner gate labeled by the operator and with
children the operand UUIDs. Similarly, data values that are results of
aggregate queries are represented as gates in the same circuit (with
UUIDv5 representation) encoding their semimodule construction.

Storage of this provenance circuit in the
DBMS is a challenge:
\begin{asparaenum}[(i)]
	\item \label{item:append-only} It is an append-only data structure; except in cases where one
	wants to perform some clean-up, there is never the need of removing
	gates from the circuit, and a gate never needs to be updated.
	\item \label{item:updated-query-time} It is a data structure which is updated at the time a query is
	evaluated, following the query rewriting approach of
	Section~\ref{sec:query_rewriting}.
	\item \label{item:persistent} The content of this circuit needs to be stored in a persistent
	manner: UUIDs become uninterpretable without the circuit.
	\item \label{item:large} The circuit can become very large, as it
	contains one gate per tuple, and one gate
	per operation performed in a query.
	\item \label{item:fast} Access to the circuit needs to be as fast as
	possible.
\end{asparaenum}

Initial implementations of ProvSQL stored this circuit as an extra table
within the database, which
solved~(\ref{item:persistent})--(\ref{item:large}). But PostgreSQL
tables are not optimized for append-only operations
(\ref{item:append-only}) and, most importantly,
(\ref{item:updated-query-time}) caused concurrency control issues. We
then experimented with shared-memory storage of the circuit, which solved
most issues but did not scale (\ref{item:large}) and made persistence
(\ref{item:persistent}) an issue in case of failure.

Our final
implementation provides a satisfactory solution to (\ref{item:append-only})--(\ref{item:fast}):
The circuit is stored outside of the database, in \emph{append-only
	files} that
are properly indexed and \emph{memory-mapped} so that the operating
system can keep most
often used fragments in RAM buffers. To avoid concurrency issues, access to
these files is managed by a \emph{single PostgreSQL worker process},
communicating with other backends through
inter-process communication (with \emph{pipes}). Finally, each backend
process has a \emph{small local cache} of most recently accessed circuit information.

\subsection{Query rewriting}

The main way ProvSQL is keeping track of the provenance of data is by the
query rewriting approach discussed in
Section~\ref{sec:query_rewriting}. To implement this within PostgreSQL, ProvSQL
defines
a custom hook at \emph{query planning time}
before the query is sent to PostgreSQL's planner. If the query involves
annotated relations (i.e., tables with the special \emph{provsql}
attribute of UUID type), the query is analyzed to determine if it is part of the
supported language (the extended relational algebra described
in Section~\ref{sec:prelim}, as represented in SQL). If not, an error message is
displayed instead of performing inadequate provenance computation. If so,
the query is rewritten as described in Section~\ref{sec:query_rewriting} and
then sent to the regular planner for continued processing and evaluation.

\subsection{Specialization to specific (m-)semirings}

Annotations computed by ProvSQL are abstract UUIDs, with description in
the circuit on how these are computed from base tuple UUIDs. One
important point is that representation can be specialized to any
arbitrary \mbox{(m-)semi}\-ring and semimodule by applying homomorphisms: ProvSQL
essentially works in the \emph{universal} semiring (the how-semiring of
integer polynomials, see~\cite{green2007provenance}) and in the
\emph{universal} m-semiring (the free m-semiring, see
\cite{dblp:journals/japll/geertsp10}), from which there exist unique
homomorphisms to any application (m-)semiring.
Each algebraic structure of interest is compiled into a C++ class
whose methods describe the different operations
($\oplus$, $\otimes$, possibly $\ominus$ and~$\delta$)
of the algebraic
structure.
At a user's request,
given a \emph{provenance
	mapping} from base UUIDs to values in that algebraic structure (e.g.,
  integers in the counting semiring),
the relevant part of the circuit is
retrieved from its memory mapped representation, and ProvSQL
  evaluates the circuit in a bottom-up
manner to determine the actual semiring annotation of every tuple in the
result of a query, using the methods provided by the class.

\subsection{Probability evaluation}
\label{sec:impl_prob}

Probabilistic query evaluation in ProvSQL relies on
Corollary~\ref{cor:pqe}: provenance is computed in the same way as for
provenance tracking, using query rewriting, and when a marginal
probability computation is requested, we specialize the provenance to
$\mathcal{B}[X]$ (given a mapping of base UUIDs to some Boolean function,
e.g., one distinct variable per UUIDs to obtain a tuple-independent
database) and then compute the probability of the corresponding Boolean
circuit. Note that the computation of Boolean provenance, as it is
crucial for probabilistic query evaluation, is further optimized from the
regular retrieval of subcircuits from the the memory-mapped circuit.

ProvSQL implements different ways of computing the probability of a
Boolean circuit (e.g., naïve enumeration of possible worlds,
Monte-Carlo sampling, or some other approximation techniques such as
WeightMC~\cite{DBLP:conf/aaai/ChakrabortyFMSV14}). We focus here on its
default computation method, which works in three steps:
\begin{asparaenum}[(1)]
	\item We determine whether the Boolean circuit is
	read-once~\cite{DBLP:journals/pvldb/SenDG10}, i.e., if every variable
	occurs only once; if so, its probability can be computed in
	linear-time.
	\item Otherwise, we use a heuristic
	algorithm~\cite{maniu2019experimental} to quickly obtain a
	tree-decomposition of small width ($\leq 10$) of the circuit, if
	possible; if so,  we use an algorithm
	from~\cite{dblp:journals/mst/amarillicms20} to extract a
	deterministic and decomposable circuit from the Boolean circuit, on
	which the probability is computed in linear time.
	\item Otherwise, we encode the circuit in conjunctive normal form using
	Tseitin's transformation \cite{Tseitin1983} and use an external
	knowledge compiler (by default d4
	\cite{dblp:conf/ijcai/lagniezm17} but we also support c2d
	\cite{DBLP:conf/ecai/Darwiche04} and DSHARP
	\cite{DBLP:conf/ai/MuiseMBH12}) as a black box to compile it into a
	deterministic and decomposable circuit, on which the probability is
	computed in linear time (but this new circuit may be exponential in
	size).
\end{asparaenum}

Recall that probabilistic query evaluation is intractable
($\mathsf{\#P}$-hard):
the potential exponential blow-up of the last step is
unavoidable.

\subsection{Other features}
At the time of writing, ProvSQL also implements other features,
which
are outside of the scope of this paper: it supports provenance for
non-terminal aggregate queries, following Section~4
of~\cite{amsterdamer2011provenance}; Shapley value
computation and expected Shapley value computation as described
in~\cite{DBLP:journals/pacmmod/KarmakarMSB24}; provenance of update
operations~\cite{widiaatmaja2025demonstration} as proposed in~\cite{DBLP:conf/sigmod/BourhisDM20}; expected
value computation for \mintinline{SQL}{COUNT}, \mintinline{SQL}{MIN},
\mintinline{SQL}{MAX}, \mintinline{SQL}{SUM}, following the algorithms
in~\cite{abiteboul2011capturing}; and where-provenance
\cite{bunemankt01,senellart2018provsql}, which is a form of provenance not captured by
semirings \cite{DBLP:journals/ftdb/CheneyCT09}.
